% !PRINTED 404
\weberpage{404}

%<<CONT
all have the same norm; they are associated according to $[Q_1]$, but
not according to $[Q]$. Accordingly we obtain from (6):

\[ \mathrm{Lim}_{s=1}\sum\frac{s-1}{(Ny_1)^s} = \frac{r}{a\sqrt D}\,\frac{\log\varepsilon}{\lambda}; \]
in order to obtain from this the sum relating to the numbers (9),
one has, since $N(\varepsilon_1)=1$, $N(r)=r^2$, to multiply by $\lambda$
and to divide by $r^2$. One obtains in this way:

\[ \mathrm{Lim}\sum\frac{s-1}{(Ny)^s} = \frac{1}{r\,a\sqrt D}\log\varepsilon. \]
If one subtracts this from the sum (6) and proceeds in the same way
with all prime factors $r$ of $Q$, there results:

\weq{10}{\mathrm{Lim}\sum\frac{s-1}{(Ny)^s} = \prod\Bigl(1-\frac1r\Bigr)\frac{1}{a\sqrt D}\log\varepsilon, \qquad D>0,}
and still more simply for a negative discriminant:

\[ = \prod\Bigl(1-\frac1r\Bigr)\frac{\pi}{a\sqrt{-D}}, \qquad D<0, \]
in which the sum extends only over those numbers $y$ that are
coprime to $Q$, and the product $\Pi$ relates to all
prime factors of $r$. % [sic] The German source prints r; the immediately preceding construction requires the prime factors of Q.
Only in the two exceptional cases
$D=-4$, $D=-3$ is the right-hand side of the last formula
still to be divided by 2 or by 3.

If one decomposes $y$ into ideal factors

\weq{11}{y = \mathfrak{a}\,\mathfrak{m},}
then $\mathfrak{m}$ runs through all ideals of a definite ideal class according to
the order $[Q]$ that are coprime to $Q$. This class is,
if $\mathfrak{a}$ belongs to the class $A$,

\[ M = A^{-1}, \]
and among the numbers

\[ N(y) = a\,N(\mathfrak{m}) \]
in the case of the positive discriminant each $N(\mathfrak{m})$ occurs only once,
in the case of a negative discriminant twice or four or
six times.

Accordingly there results from (10):

\weq{12}{\begin{aligned}
\mathrm{Lim}_{s=1}\sum\frac{s-1}{(N\mathfrak{m})^s} &= \prod\Bigl(1-\frac1r\Bigr)\frac{1}{\sqrt D}\log\varepsilon, & D&>0\\
&= \prod\Bigl(1-\frac1r\Bigr)\frac{\pi}{\sqrt{-D}}, & D&<0
\end{aligned}}
%CONT>>
